\section{The scroll bound for a fixed shift}
\label{app:scroll-shift}

We refine the fixed-shift bound of Theorem~\ref{thm:translation} using
the full set of linear syzygies of a rational normal scroll.
Castryck--Cools--Demeyer--Lemmens
use the linear strand of an ambient scroll for a toric surface
\cite[Sec.~4.1]{CCDL19}. Here the Goldbach grading records the weight
of each scroll generator.

\begin{proposition}[The scroll bound for a fixed shift]\label{thm:scroll-shift}
Fix $h>0$, let $B\subseteq\A\cap(\A-h)$ be finite, and put $r=|B|$.
For the binary Goldbach ring on $\A$, over every field and for $i\ge1$,
\[
 \beta_{i,n,i+1}\ \ge\
 \sum_{j=1}^{i}\#\left\{\sigma\subseteq B:\
       |\sigma|=i+1,\quad \sum_{b\in\sigma} b+jh=n\right\}.
\]
In particular $\beta_{i,*,i+1}\ge i\binom{r}{i+1}$.
\end{proposition}
\begin{proof}
We may assume $r\ge i+1$. In the polynomial ring on the labels of
$B\cup(B+h)$, form the matrix $M$ with columns $(x_b,x_{b+h})^{\mathsf T}$
for $b\in B$. Its minors belong to $I$, since $b+(b'+h)=(b+h)+b'$.
The edges $b\to b+h$ form disjoint paths. Each vertex has at most one
predecessor and one successor, and labels increase along each edge.
Thus the columns of a path with vertices $a_0,\ldots,a_s$, in order,
form a block
\[
 \begin{pmatrix}x_{a_0}&\cdots&x_{a_{s-1}}\\
                 x_{a_1}&\cdots&x_{a_s}\end{pmatrix}.
\]
Different blocks use disjoint variables. This is a rational normal
scroll matrix. If there are $t$ paths, there are $r+t$ variables. Over
an algebraic closure, on the locus of the minor equations where the top row
is nonzero, the bottom row
equals a scalar multiple of it. That scalar and the initial variable
of each path parametrize this locus, of dimension $t+1$. Where the
top row vanishes only the $t$ terminal variables are free. Hence the
minor ideal has height $r-1$, and the Eagon--Northcott complex is a
minimal resolution of its quotient \cite{EN62}. Extending it to $R$ by the remaining
variables preserves exactness and minimality, also for infinite $\A$.

To read its bidegrees, assign weight $b$ and monomial degree one to the column
indexed by $b$, and give the two row basis elements weights $0,-h$
and monomial degree zero. This makes $M$ homogeneous. In homological degree $i$ an
Eagon--Northcott generator uses $i+1$ columns indexed by $\sigma$,
the dual row determinant, and a divided power of degree $i-1$ in
the dual rows. The determinant contributes weight $h$. Taking $j-1$
copies of the second dual row contributes another $(j-1)h$, for
$j=1,\ldots,i$. Thus the generators have precisely the weights and
monomial degrees in the formula.

The quotient map $R/I_2(M)\to G$ induces an injection on the linear
strand. Indeed, in monomial degree $i+1$, the Koszul terms in homological degrees
$i+1$ and $i$ have coefficient monomial degrees zero and one. The quotient map
identifies both terms and their incoming differential. It takes each
cycle to a cycle, and both complexes have the same boundaries in
homological degree $i$. This proves the weightwise bound. Summing over $n$
counts each $(i+1)$-subset $i$ times.
\end{proof}

The classes $z_{\sigma,h}$ give the $j=1$ contribution. To see this,
first use independent entries $X_b,Y_b$ in the two rows of $M$.
Give both entries their column degree and count each lower-row entry
once in a further grading. The chain \eqref{eq:translated-cycle}, with
coefficient $Y_b$ and exterior generator indexed by $X_b$, is a cycle
modulo the minors. Its column multidegree is $\sigma$ and its lower-row
degree is one. The scroll resolution has just one generator in this
fine degree. Substituting $x_b,x_{b+h}$ for $X_b,Y_b$ gives
$z_{\sigma,h}$, whose image in $G$ is nonzero by
Theorem~\ref{thm:translation}.

For example, $B=\{3,11,17\}$ and $h=2$ give two cubic classes, at
weights $33$ and $35$. The first is the introductory cycle. The second
can be written
\[
 x_3\varepsilon_{13}\varepsilon_{19}
 -x_{11}\varepsilon_5\varepsilon_{19}
 +x_{17}\varepsilon_5\varepsilon_{13}.
\]
It comes from interchanging the rows of $M$. The proposition counts
these classes for one shift at a time. For instance the shifts
$h=4$ on $\{3,13\}$ and $h=10$ on $\{3,7\}$ give cycles whose
difference is $\mathrm{d}(\varepsilon_7\varepsilon_{13})$.
Thus the fixed-shift bounds need not add across different shifts.
